\begin{solapartat}
\punts{0,25} A partir del gràfic $2T = 17$ dies $\Rightarrow T = 8,5$ dies $= 7,34\times 10^{5}\ \mathrm{s}$.

De la tercera llei de Kepler i tenint en compte que la constant de Kepler és igual a la del nostre sistema solar:

\punts{0,7}
\[ \frac{r_\mathrm{orb\ terra}^3}{T_\mathrm{terra}^2} = \frac{r_\mathrm{orb\ planeta}^3}{T_\mathrm{planeta}^2} \;\Rightarrow\; r_\mathrm{orb\ planeta} = r_\mathrm{orb\ terra}\left(\frac{T_\mathrm{planeta}}{T_\mathrm{terra}}\right)^{2/3} = 1,22\times 10^{10}\ \mathrm{m} = 8,15\times 10^{-2}\ \mathrm{ua} \]
\punts{0,3}

Cal tenir present que aquest no és l'únic plantejament correcte; per exemple, també es pot resoldre a partir de la segona llei de Newton i la llei de la gravitació universal de Newton.
\end{solapartat}

\begin{solapartat}
\punts{0,45} $v_\mathrm{planeta} = \dfrac{2\pi r_\mathrm{orb\ planeta}}{T} = 1,05\times 10^{5}\ \mathrm{m/s}$ \punts{0,2}

\punts{0,4} $a_n = \dfrac{v^2}{r} = \dfrac{(1,05\times 10^{5})^2}{1,22\times 10^{10}} = 0,895\ \mathrm{m/s^2}$ \punts{0,2}
\end{solapartat}
